\documentclass{article}
\usepackage[utf8]{inputenc}
\usepackage[left=0.5in,right=0.5in,top=2cm,bottom=2cm]{geometry}
\usepackage{crop,graphicx,amsmath,array,color,amssymb,flushend,stfloats,amsthm,chngpage,times,fancyhdr,lipsum,lastpage,subcaption,cancel,algpseudocode,hyperref}

%%%%%%%%%%%%   Header and Footer  %%%%%%%%%%%%%
\pagestyle{fancy}

\fancypagestyle{plain}{%
  \renewcommand{\headrulewidth}{0pt}%
  \fancyhf{}%
  \fancyfoot[R]{Page \bf\thepage\ \rm of \bf\pageref{LastPage}}%
}


%%%% Customise Titles and Headers: %%%%
\title{Bug Report Classifier - Replication}
\author{Jack Honour [3063159]}
\date{}

\fancyhf{}
\fancyfoot[L]{Jack Honour [3063159]}
\fancyhead[R]{University of Birmingham \LaTeX{}}
\fancyfoot[R]{Page \bf\thepage\ \rm of \bf\pageref{LastPage}}
\fancyhead[L]{Bug Report Classifier - Replication}


\begin{document}

%%%%%%%%%%%% Make Title and Format Lines %%%%%%%%%%%%
\maketitle											%
\vspace{-112px}										%
\noindent\rule{\linewidth}{1pt} \par				%
\vspace{100px}										%
\vspace{-20px}										%
\noindent\rule{\linewidth}{1pt} \par				%
\vspace{10px}										%
%%%%%%%%%%%%%%%%%%% Insert Logos %%%%%%%%%%%%%%%%%%%%	
\vspace{-85px}										%
\noindent											%
\begin{minipage}{0.5\textwidth}\begin{flushleft}	%
\hspace{20px}										%
\includegraphics[scale = 0.06]{Resources/UoB}		%
\end{flushleft}\end{minipage}						%
\begin{minipage}{0.5\textwidth}\begin{flushright}	%
\includegraphics[scale = 0.06]{Resources/UoB}		%
\hspace*{20px}										%
\end{flushright}\end{minipage}						%
\vspace{20px}										%
%%%%%%%%%%%%%%%% Headers and Footers %%%%%%%%%%%%%%%%

\fontsize{11}{16}\selectfont
\setlength{\parindent}{0pt} % Paragraph indent
\setlength{\parskip}{4pt}   % Space between paragraphs



\section*{IMPORTANT}
You MUST use git to clone the repository and run the code, otherwise you will not be able to replicate the results. Do NOT download the code as a zip file from GitHub, as this will not include the necessary LFS files which contain the pre-trained models.

\begin{verbatim}
git lfs install
git clone https://github.com/jaxkdev/ise_coursework
\end{verbatim}

\textit{Note: cloning will take as long as it takes to download 2.5GB of data, so please be patient.}

Please read the \texttt{requirements.pdf} file before running anything to ensure you have the necessary hardware and software requirements. As well as a correct virtual environment setup with the required dependencies installed.

\subsection*{Data Accuracy}
The results will NOT be the exact same as the original results due to internal parallelism in the algorithms and the use of internal random seeds.
However, where possible I have used the same seed and the results should be very close to the original results, with similar means and standard deviations across all metrics.

Please please dont mark down for not being exactly the same, as this is not possible I tried very hard,
even going down to 1 process ($n\_jobs=1$) to try but still small variations on different devices.

\vspace{10px}

\section*{Replication}

\begin{verbatim}
(venv) python src/dataset.py  # To view dataset details.
(venv) python analysis.py     # To run full replication of all results.
\end{verbatim}

This will run the full analysis, replicating the results for all algorithms across all projects. Note that the improved\_6 algorithm (BERT-based) is very time-consuming to train, so by default it will only run one repetition for each project.

The replicated results will be saved in the \texttt{results\_replicated} folder, with the same structure as the original \texttt{results} folder. Each algorithm will have its own subfolder, and within each subfolder there will be CSV files for each project and repetition.

\section*{Custom Replication}

The following arguments can be passed to \texttt{analysis.py} to customise the replication process:
\begin{itemize}
    \item \texttt{--algorithm}: Specify which algorithm(s) to replicate. Can be one or more of \texttt{baseline}, \texttt{improved\_1}, \texttt{improved\_2}, \texttt{improved\_3}, \texttt{improved\_4}, \texttt{improved\_5}, \texttt{improved\_6}, \texttt{final}, \texttt{*}. Default is * algorithms.
    \item \texttt{--project}: Specify which project(s) to replicate. Can be one or more of \texttt{caffe}, \texttt{incubator-mxnet}, \texttt{keras}, \texttt{pytorch}, \texttt{tensorflow}, \texttt{all}, \texttt{*}. Default is * projects.
    \item \texttt{--repetitions}: Specify how many repetitions to run for each algorithm and project. Default is 10, 20 and 50 repetitions combined (except for improved\_6 which defaults to 1 repetition due to time constraints).
    \item \texttt{--use-cached-models}: Whether to use cached pre-trained models if available, or always train from scratch. Default is to use cached models if available. Can be set to \texttt{true} or \texttt{false}.
\end{itemize}

Example, to replicate only the final algorithm all individual projects with 2 repetitions, you would run:
\begin{verbatim}
(venv) python analysis.py --algorithm=final --project=* --repetitions=2
\end{verbatim}

Not to be confused with the all-projects analysis, which runs each algorithm on a combined dataset of all projects. To replicate the all-projects analysis for the final algorithm with 5 repetitions, you would run:
\begin{verbatim}
(venv) python analysis.py --algorithm=final --project=all --repetitions=5
\end{verbatim}


\section*{Tool Usage Replication}

\begin{verbatim}
(venv) python tool.py --file=test.txt
(venv) python tool.py --file=test.txt --algorithm=baseline
(venv) python tool.py --file=test_pos.txt
(venv) python tool.py --file=test_pos.txt --algorithm=baseline
\end{verbatim}

These commands will run the tool on the provided sample file, using the default final algorithm or the baseline algorithm as specified. You can replace \texttt{test.txt} and \texttt{test\_pos.txt} with any text file containing bug reports to classify them using the tool.

\end{document}